%% toc: P15 165 165.6
\section{Partition updates and conglomerability}
\label{sec:annex-p15-w02b-p15-034-partition-updates}

Throughout, \(\underline P\) is a coherent lower prevision on
\(\mathcal L(\Omega)\). By \cref{thm:p15-ip-linear-domain} and
\cref{thm:p15-ip-natural-extension} it satisfies
\(\inf h\leq\underline P(h)\leq\sup h\), positive homogeneity,
superadditivity, monotonicity and
\(|\underline P(h)-\underline P(k)|\leq\|h-k\|_\infty\). Two consequences
are used repeatedly. First, \(\underline P(0)=0\) and
\(\underline P(B)\geq\inf I_B\geq0\) for every event. Second, for all
gambles \(u,v\),
\[
 \sup_{\Omega}(u-v)\geq\underline P(u)-\underline P(v),\tag{\(*\)}
\]
because superadditivity gives
\(\underline P(v)\geq\underline P(u)+\underline P(v-u)
\geq\underline P(u)+\inf(v-u)=\underline P(u)-\sup(u-v)\).

\begin{annexitem}{annex:p15.w02b-p15-034.cell-facts}{thm:p15-w02b-p15-034-cell-facts}{Cellwise facts}
\begin{proof}
(1) Applying (C2) with \(\lambda=2\) to \(f=0\) gives
\(\underline P(0\mid B)=2\underline P(0\mid B)\), so
\(\underline P(0\mid B)=0\). By (C3),
\(\underline P(f\mid B)+\underline P(-f\mid B)\leq\underline P(0\mid B)=0\),
so by (C1)
\(\underline P(f\mid B)\leq-\underline P(-f\mid B)\leq-\inf_B(-f)=\sup_Bf\).
The lower bound is (C1). If \(f=g\) on \(B\), then (C3) and (C1) give
\(\underline P(f\mid B)\geq\underline P(g\mid B)+\underline P(f-g\mid B)
\geq\underline P(g\mid B)+\inf_B(f-g)=\underline P(g\mid B)\), and the
reverse inequality follows by exchanging \(f\) and \(g\). The bounds give
\(|\underline P(f\mid B_\omega)|\leq\|f\|_\infty\) for every \(\omega\), so
\(\underline P(f\mid\mathcal B)\) is bounded. Finally, on \(B\) the gambles
\(I_Bf\) and \(f\) agree, so \(G(I_Bf\mid\mathcal B)=f-\underline P(f\mid B)\)
there; on a cell \(B'\ne B\), \(I_Bf\) vanishes, so
\(\underline P(I_Bf\mid B')=\underline P(0\mid B')=0\) and
\(G(I_Bf\mid\mathcal B)=0\) there. Hence
\(G(I_Bf\mid\mathcal B)=I_B(f-\underline P(f\mid B))=G(f\mid B)\).

(2) By (C3) and (C1),
\(\underline P(g\mid B)\geq\underline P(f\mid B)+\underline P(g-f\mid B)
\geq\underline P(f\mid B)+\inf_B(g-f)\). Thus
\(\sup_B(f-g)=-\inf_B(g-f)\geq\underline P(f\mid B)-\underline P(g\mid B)\),
which is the claim.

(3) For real \(q,q'\),
\(\|I_B[f-q]-I_B[f-q']\|_\infty\leq|q-q'|\), so the Lipschitz property of
\(\underline P\) gives \(|F(q)-F(q')|\leq|q-q'|\). For \(\delta>0\),
\(I_B[f-q]=I_B[f-q-\delta]+\delta I_B\); superadditivity and positive
homogeneity give \(F(q)\geq F(q+\delta)+\delta\underline P(B)\). Since
\(\underline P(B)\geq0\), \(F\) is nonincreasing, and if
\(\underline P(B)>0\) it is strictly decreasing, so it has at most one
root. In that case \(F(\inf_Bf)=\underline P(I_B[f-\inf_Bf])\geq0\) because
the gamble is nonnegative, and \(F(\sup_Bf)\leq\sup I_B[f-\sup_Bf]\leq0\).
By continuity and the intermediate value theorem a root exists in
\([\inf_Bf,\sup_Bf]\).

(4) Let \(B_1,\ldots,B_k\) be distinct cells with
\(\underline P(B_i)\geq1/m\). They are disjoint, so
\(I_{B_1\cup\cdots\cup B_k}=\sum_iI_{B_i}\leq1\). Superadditivity and the
upper bound give
\(1\geq\underline P(I_{B_1\cup\cdots\cup B_k})\geq\sum_i\underline P(B_i)\geq k/m\),
so \(k\leq m\). The set of cells of positive lower probability is the union
over \(m\geq1\) of the sets \(\{B:\underline P(B)\geq1/m\}\), each of which
has at most \(m\) members; it is finite or countably infinite.

(5) For disjoint \(B_1,\ldots,B_k\),
\(I_{B_1\cup\cdots\cup B_k}f=\sum_iI_{B_i}f\), and superadditivity gives
\(\underline P(I_{B_1\cup\cdots\cup B_k}f)\geq\sum_i\underline P(I_{B_i}f)\geq0\).
If only finitely many cells have positive lower probability, there is no
infinite sequence of distinct such cells, and the condition of
\cref{def:p15-w02b-p15-034-conglomerable} holds vacuously. A finite partition
has finitely many cells; a partition whose cells all have lower probability
zero has none of positive lower probability.
\end{proof}
\end{annexitem}

\begin{annexitem}{annex:p15.w02b-p15-034.necessity}{thm:p15-w02b-p15-034-necessity}{Two portfolios and one infinite portfolio}
\begin{proof}
(1) Fix \(f\) and \(B\), and put \(h=G(f\mid B)=I_B(f-\underline P(f\mid B))\).
By \cref{thm:p15-w02b-p15-034-cell-facts}(1),
\(G(I_Bf\mid\mathcal B)=h\). Apply (JC1) with \(f_1=0\), \(g_1=I_Bf\) and
\(f_2=h\). Since \(G(0)=0-\underline P(0)=0\), the net gain is
\(h-(h-\underline P(h))=\underline P(h)\), a constant, so (JC1) gives
\(\underline P(h)\geq0\). Apply (JC2) with \(f_1=h\), \(g_1=0\),
\(g_2=f\) and \(B_0=B\). Since \(G(0\mid\mathcal B)=0\) by
\cref{thm:p15-w02b-p15-034-cell-facts}(1), the net gain is
\(Z=(h-\underline P(h))-h=-\underline P(h)\), again a constant. Every set
in the prescribed family is nonempty, so its supremum over any of them is
\(-\underline P(h)\), and (JC2) gives \(\underline P(h)\leq0\). Hence
\(\underline P(h)=0\).

(2) If \(\underline P(B)>0\), item (1) says that
\(\underline P(f\mid B)\) is a root of \(q\mapsto\underline P(I_B[f-q])\),
and \cref{thm:p15-w02b-p15-034-cell-facts}(3) says that this root is
unique. If \(\underline P(B)=0\), the bound is (C1).

(3) Let \(B_1,B_2,\ldots\) be distinct cells with
\(\underline P(B_n)>0\) and let \(\underline P(I_{B_n}f)\geq0\) for all
\(n\). Put \(F_n(q)=\underline P(I_{B_n}[f-q])\) and
\(r_n=\underline P(f\mid B_n)\). By (1), \(F_n(r_n)=0\). If \(r_n<0\),
\cref{thm:p15-w02b-p15-034-cell-facts}(3) would give
\(F_n(0)\leq F_n(r_n)-(0-r_n)\underline P(B_n)<0\), contradicting
\(F_n(0)=\underline P(I_{B_n}f)\geq0\). So \(r_n\geq0\) for all \(n\).

Let \(D=\bigcup_nB_n\) and \(f'=I_Df\). For each \(n\), \(f'=f\) on
\(B_n\), so \(\underline P(f'\mid B_n)=r_n\geq0\) by
\cref{thm:p15-w02b-p15-034-cell-facts}(1). A cell \(B\) not among the
\(B_n\) is disjoint from \(D\), so \(f'=0\) on \(B\) and
\(\underline P(f'\mid B)=\underline P(0\mid B)=0\). Hence the gamble
\(\underline P(f'\mid\mathcal B)\) is nonnegative everywhere. Apply (JC1)
with \(f_1=0\), \(g_1=f'\), \(f_2=f'\). The net gain is
\[
 G(f'\mid\mathcal B)-G(f')
 =\bigl(f'-\underline P(f'\mid\mathcal B)\bigr)-\bigl(f'-\underline P(f')\bigr)
 =\underline P(f')-\underline P(f'\mid\mathcal B),
\]
whose supremum over \(\Omega\) is
\(\underline P(f')-\inf_\omega\underline P(f'\mid B_\omega)\). (JC1) makes it
nonnegative, so
\(\underline P(I_Df)=\underline P(f')\geq\inf_\omega\underline P(f'\mid B_\omega)\geq0\).
This is the conglomerability inequality. The portfolio used is active on
every cell \(B_n\) at once.
\end{proof}
\end{annexitem}

\begin{annexitem}{annex:p15.w02b-p15-034.existence}{thm:p15-w02b-p15-034-existence}{Constructing the least update under conglomerability}
\begin{proof}
If a partition update exists, \(\underline P\) is
\(\mathcal B\)-conglomerable by
\cref{thm:p15-w02b-p15-034-necessity}(3).

Conversely assume \(\mathcal B\)-conglomerability.
\emph{Well defined.} For a cell with \(\underline P(B)>0\), the root
exists and is unique by \cref{thm:p15-w02b-p15-034-cell-facts}(3).

\emph{Separate coherence.} On a cell of lower probability zero,
\(\inf_B\) satisfies (C1) with equality, (C2) because \(\lambda>0\), and
(C3) because \(\inf_B(f+g)\geq\inf_Bf+\inf_Bg\). Let
\(\underline P(B)>0\), write \(F_f(q)=\underline P(I_B[f-q])\) and let
\(q_f\) be its root. For (C1), suppose \(q_f<\inf_Bf\); the strict
decrease in \cref{thm:p15-w02b-p15-034-cell-facts}(3) would give
\(F_f(\inf_Bf)\leq F_f(q_f)-(\inf_Bf-q_f)\underline P(B)<0\), whereas
\(F_f(\inf_Bf)\geq0\) because \(I_B[f-\inf_Bf]\geq0\). For (C2),
\(F_{\lambda f}(\lambda q_f)=\lambda F_f(q_f)=0\), so
\(q_{\lambda f}=\lambda q_f\) by uniqueness. For (C3), superadditivity gives
\(F_{f+g}(q_f+q_g)\geq F_f(q_f)+F_g(q_g)=0\); if \(q_{f+g}<q_f+q_g\),
strict decrease would give \(F_{f+g}(q_f+q_g)<F_{f+g}(q_{f+g})=0\). Hence
\(q_{f+g}\geq q_f+q_g\).

\emph{Every conditional gain is almost acceptable.} Fix \(g\) and put
\(G=G(g\mid\mathcal B)\), computed with \(\underline P_{\mathcal B}\). By
\cref{thm:p15-w02b-p15-034-cell-facts}(4) the cells of positive lower
probability can be listed without repetition as \(B_1,B_2,\ldots\), a finite
or infinite list; let \(D\) be their union (empty if there are none). For
each listed cell, \(I_{B_n}G=G(g\mid B_n)\) and
\(\underline P(I_{B_n}G)=F_g(q_g)=0\) at that cell. If the list is
infinite, \(\mathcal B\)-conglomerability applied to the gamble \(G\) gives
\(\underline P(I_DG)\geq0\). If it is finite and nonempty,
\cref{thm:p15-w02b-p15-034-cell-facts}(5) gives the same inequality; if it
is empty, \(I_DG=0\). On a cell \(B\) of lower probability zero,
\(G=g-\inf_Bg\geq0\). Therefore \(G\geq I_DG\) pointwise and, by
monotonicity,
\[
                  \underline P(G(g\mid\mathcal B))\geq\underline P(I_DG)\geq0.
\]

\emph{Each cell's own gain has lower prevision zero.} Let
\(k=G(g_2\mid B_0)\). If \(\underline P(B_0)>0\),
\(\underline P(k)=0\) by the definition of the root. If
\(\underline P(B_0)=0\), then \(k=I_{B_0}(g_2-\inf_{B_0}g_2)\), so
\(0\leq k\leq cI_{B_0}\) with \(c=\sup k\geq0\); monotonicity and
homogeneity give \(0\leq\underline P(k)\leq c\,\underline P(B_0)=0\).

\emph{(JC1).} Put \(u=f_1+G(g_1\mid\mathcal B)\). The net gain is
\(u-f_2-\underline P(f_1)+\underline P(f_2)\). By \((*)\) its supremum over
\(\Omega\) is at least
\(\underline P(u)-\underline P(f_1)
\geq\underline P(f_1)+\underline P(G(g_1\mid\mathcal B))-\underline P(f_1)\geq0\).

\emph{(JC2).} Let \(k=G(g_2\mid B_0)\) as above, so
\(\underline P(k)=0\). If \(f_1\ne0\), \(\Omega\) belongs to the prescribed
family. With \(u\) as before, \(Z=u-k-\underline P(f_1)\), and by \((*)\)
\(\sup_\Omega Z\geq\underline P(u)-\underline P(k)-\underline P(f_1)
\geq\underline P(G(g_1\mid\mathcal B))\geq0\). If \(f_1=0\), then
\(G(f_1)=0\) and \(B_0\) belongs to the family. On \(B_0\),
\(Z=(g_1-\underline P_{\mathcal B}(g_1\mid B_0))-(g_2-\underline P_{\mathcal B}(g_2\mid B_0))\),
and \cref{thm:p15-w02b-p15-034-cell-facts}(2) gives
\(\sup_{B_0}Z\geq0\).

So \(\underline P_{\mathcal B}\) is a partition update. For any partition
update, \cref{thm:p15-w02b-p15-034-necessity}(2) forces the root on cells of
positive lower probability and the bound \(\geq\inf_Bf\) on the others,
which are the displayed relations.
\end{proof}
\end{annexitem}

\begin{annexitem}{annex:p15.w02b-p15-034.counterexample}{thm:p15-w02b-p15-034-counterexample}{The finitely additive mixture computed in full}
\begin{proof}
(1) On \(\mathcal L(\mathbb N_+)\) let
\(Q(h)=\liminf_{n\to\infty}h(n)\). Then \(Q(h)\geq\inf h\),
\(Q(th)=tQ(h)\) for \(t>0\), and
\(\liminf(h+k)\geq\liminf h+\liminf k\). By
\cref{thm:p15-ip-linear-domain}, applied with the power-set algebra of
\(\mathbb N_+\), \(Q\) is coherent; in particular it avoids sure loss, so by
\cref{thm:p15-ip-representation} its compatible set is nonempty. Choose
\(\nu\) in it: \(\nu\) is positive, normalized and linear, and
\(\nu(h)\geq\liminf h\). For finite \(F\), \(\nu(I_F)\geq\liminf I_F=0\),
and
\(\nu(I_F)=1-\nu(I_{\mathbb N_+\setminus F})\leq1-\liminf I_{\mathbb N_+\setminus F}=0\).

(2) The series defining \(\mu\) converges absolutely, and
\[
 |\mu(f)|\leq\|f\|_\infty\sum_n2^{-n}=\|f\|_\infty;
\]
\(\mu\) is linear,
positive and \(\mu(1)=1\). The map \(f\mapsto(n\mapsto f(n,1))\) is linear,
positive and sends \(1\) to \(1\), so \(\nu_1\) has the same three
properties, and so does their average \(P\). For such a functional
\(P(g)-\inf g=P(g-\inf g)\geq0\), positive homogeneity holds and
superadditivity holds with equality, so \(P\) is coherent by
\cref{thm:p15-ip-linear-domain}; its conjugate is
\(-P(-g)=P(g)\).

Now \(\mu(I_{B_n})=2^{-n}\) and \(\nu_1(I_{B_n})=\nu(I_{\{n\}})=0\), so
\(P(B_n)=2^{-n-1}\) and \(\sum_nP(B_n)=1/2\). Next
\(\mu(I_{Y_1})=\sum_n2^{-n}/2=1/2\) and \(\nu_1(I_{Y_1})=\nu(1)=1\), so
\(P(Y_1)=3/4\). Finally
\(P(I_{B_n}I_{Y_1})=\tfrac12\bigl(2^{-n}/2+\nu(I_{\{n\}})\bigr)=2^{-n-2}\),
whose ratio to \(P(B_n)\) is \(1/2\).

(3) Here \(f(n,0)=1\) and \(f(n,1)=-1\). Then
\(\mu(I_{B_n}f)=2^{-n}(1-1)/2=0\) and
\(\nu_1(I_{B_n}f)=\nu(-I_{\{n\}})=0\), so \(P(I_{B_n}f)=0\). Also
\(\mu(f)=0\) and \(\nu_1(f)=\nu(-1)=-1\), so \(P(f)=-1/2\). The sequence
\(B_1,B_2,\ldots\) consists of distinct cells of positive probability, each
with \(P(I_{B_n}f)\geq0\), and its union is \(\Omega\); but
\(P(I_\Omega f)=-1/2<0\). So \(P\) is not \(\mathcal B\)-conglomerable, and
by \cref{thm:p15-w02b-p15-034-existence} it has no partition update.

(4) The average \(a_n\) is linear, \(a_n(1)=1\), and
\(a_n(g)\geq\inf_{B_n}g\), so it satisfies (C1)--(C3) on \(B_n\). For real
\(q\),
\[
 P(I_{B_n}[g-q])=\tfrac12\Bigl(2^{-n}\bigl(a_n(g)-q\bigr)
 +\bigl(g(n,1)-q\bigr)\nu(I_{\{n\}})\Bigr)=2^{-n-1}\bigl(a_n(g)-q\bigr),
\]
whose only root is \(a_n(g)\). For the gamble \(f\) of item (3),
\(a_n(f)=0\) for every \(n\), so
\(G(f\mid\mathcal B)=f\), while \(G(f)=f-P(f)=f+1/2\). The net gain
\(G(f\mid\mathcal B)-G(f)\) is the constant \(-1/2\); this is the
failure of (JC1) with \(f_1=0\), \(g_1=f_2=f\).

(5) Let \(J\subseteq\mathbb N_+\) index a sequence of distinct cells and
\(D=\bigcup_{n\in J}B_n\). Since \(I_Df\) agrees with \(f\) on \(B_n\) for
\(n\in J\) and vanishes on the other cells, the defining series gives
\(\mu(I_Df)=\sum_{n\in J}2^{-n}a_n(f)=\sum_{n\in J}\mu(I_{B_n}f)\), an
absolutely convergent sum of nonnegative terms whenever each
\(\mu(I_{B_n}f)\geq0\). So \(\mu\) is \(\mathcal B\)-conglomerable. It is
coherent by the argument in (2), and \(\mu(B_n)=2^{-n}>0\). Since
\(\mu(I_{B_n}[g-q])=2^{-n}(a_n(g)-q)\), the least partition update of
\(\mu\) is \(a_n\) on \(B_n\). Finally \(\mu(Y_1)=1/2=a_n(I_{Y_1})\).
\end{proof}
\end{annexitem}

\begin{annexitem}{annex:p15.w02b-p15-034.williams}{thm:p15-w02b-p15-034-williams}{Finitely supported portfolios have expectation zero}
\begin{proof}
Every allowed conditioning event has positive probability:
\(P(\Omega)=1\) and \(P(B_n)=2^{-n-1}\). For each \(i\),
\(P\bigl(I_{C_i}(f_i-P(f_i\mid C_i))\bigr)=0\): for \(C_i=\Omega\) this is
\(P(f_i)-P(f_i)=0\) by linearity, and for \(C_i=B_n\) it is the displayed
identity in item (4) of the preceding proof at \(q=a_n(f_i)\). Let \(H\) be
the bracketed gamble. It is a finite linear combination of these gambles,
so \(P(H)=0\). Put \(U=C_0\cup\cdots\cup C_k\); \(H\) vanishes outside
\(U\). Suppose \(\sup_UH=-\epsilon<0\). Then \(H\leq-\epsilon I_U\)
pointwise, and positivity and linearity give
\(P(H)\leq-\epsilon P(U)\leq-\epsilon P(C_0)<0\), a contradiction. Hence
\(\sup_UH\geq0\).
\end{proof}
\end{annexitem}
